\documentclass[a4paper,11pt,twocolumn]{article}

\usepackage[utf8]{inputenc}
\usepackage[T1]{fontenc}
\usepackage[scaled=0.92]{helvet}
\usepackage{amssymb}
\usepackage{lmodern}
\usepackage{amsthm}
\usepackage{amsmath}
\usepackage[top=1cm, bottom=1.8cm, left=1cm, right=1cm]{geometry}
\usepackage{fancyhdr}
\usepackage[svgnames]{xcolor}
\usepackage{listings}
\usepackage{alltt}
\usepackage{draftwatermark}
\usepackage{enumerate}
\usepackage[french,vlined,boxruled]{algorithm2e}
\usepackage{graphicx}
\usepackage[francais]{babel}
\pagestyle{fancy}

\SetWatermarkText{\includegraphics{../logo.png}}
\SetWatermarkScale{0.5}
\SetWatermarkAngle{0}

%opening
\title{\vspace{-0.8cm} \textbf{Fiche de SDA} \vspace{-1cm}}
\date{}

\fancyhf{}
\renewcommand{\headrulewidth}{0pt}
\fancyfoot[C]{\tiny Che Bedara - BDE Télécom ParisTech}
\fancyfoot[RO]{\thepage}
\fancyfoot[LE]{\thepage}

\usepackage{amsthm} %ou \usepackage{ntheorem}
\usepackage{amsmath}
\usepackage{amssymb}
\usepackage{mathrsfs}
\usepackage{amsfonts}

\definecolor{vert}{rgb}{0,0.6,0}


% Une macro pour obtenir de grandes fractions dans les formules en ligne.
\def\frc#1#2{\displaystyle{#1\over#2}}

% Une macro pour les vecteurs qui donne de meilleurs résultats que \overrightarrow.
\def\vect#1{%
	\vbox{\lineskip=-.04em\baselineskip=0pt
	\halign{##\cr
	\leaders\hbox{$\scriptstyle{-}$\kern-.4em}\hfil$\scriptstyle{\rightarrow}$\cr
	$#1$\cr}}}

% Majuscules d'anglaise.
\DeclareSymbolFont{rsfscript}{U}{rsfs}{m}{n}
\DeclareSymbolFontAlphabet{\mathrsfs}{rsfscript}
\newcommand\scr{\mathrsfs}


% Des macros pour les notations usuelles.
\newcommand{\ensemblenombre}[1]{\mathbf{#1}}
\newcommand{\N}{\ensemblenombre{N}}
\newcommand{\Z}{\ensemblenombre{Z}}
\newcommand{\Q}{\ensemblenombre{Q}}
\newcommand{\R}{\ensemblenombre{R}}
\newcommand{\C}{\ensemblenombre{C}}
\newcommand{\K}{\ensemblenombre{K}}
\newcommand{\U}{\ensemblenombre{U}} % Groupe du cercle unité complexe
\newcommand{\T}{\ensemblenombre{T}} % Tore [0;2π[
\newcommand{\mes}{\ensemblenombre{M}} % Espace de mesures
\newcommand\M{\mathfrak{M}}
\newcommand\E{\mathcal{E}}
\newcommand\parties{\mathcal{P}}
\newcommand\GL{\mathcal{GL}}
\newcommand\Sym{\mathcal{S}}
\newcommand\aSym{\mathcal{A}}
\newcommand\proba{\mathbf{P}}
\newcommand\esp{\mathbf{E}}
\newcommand\Orth{\mathcal{O}}
\newcommand\cont{\mathcal{C}}
\newcommand\normale{\mathcal{N}}
\newcommand\li{[\![}
\newcommand\ri{]\!]}
\newcommand{\diff}{\mathop{}\mathopen{}\mathrm{d}}
\newcommand{\abs}[1]{\left\lvert#1\right\rvert}
\newcommand{\norme}[1]{\left\lVert#1\right\rVert}
\newcommand{\norm}[1]{\left\lVert#1\right\rVert}
\newcommand{\transp}[1]{\ensuremath{#1^\mathsf{T}}}
\newcommand{\scal}[2]{\left\langle #1 \mid #2 \right\rangle}
\newcommand{\compl}[1]{{#1}^{\mathcal{C}}} % symbole du complémentaire en exposant
\newcommand\indep{\protect\mathpalette{\protect\independenT}{\perp}} % symbole d'indépendance en probas
\def\independenT#1#2{\mathrel{\rlap{$#1#2$}\mkern3mu{#1#2}}}
\newcommand\rel{\mathcal{R}} % Pour les relations binaires
\newcommand\rar{\rightarrow}
\newcommand\lar{\leftarrow}

% Notation d'ensembles en algèbre
\newcommand\Hom{\mathrm{Hom}}
\newcommand\End{\mathrm{End}} % Endomorphismes
\newcommand\Isom{\mathrm{Isom}} % Isométries
\newcommand\Aut{\mathrm{Aut}} % Automorphismes
\newcommand\Int{\mathrm{Int}}


%
% Intervalles
\newcommand{\intervalle}[4]{\mathopen{#1}#2\mathclose{}\mathpunct{};#3\mathclose{#4}}
\newcommand{\inff}[2]{\intervalle{[}{#1}{#2}{]}}
\newcommand{\inof}[2]{\intervalle{]}{#1}{#2}{]}}
\newcommand{\info}[2]{\intervalle{[}{#1}{#2}{[}}
\newcommand{\inoo}[2]{\intervalle{]}{#1}{#2}{[}}
\newcommand{\iniff}[2]{\intervalle{[\![}{#1}{#2}{]\!]}}

% Secondes définitions d'intervalles, taille ajustable
\newcommand{\genericinterval}[4]
{
	\mathopen{}\mathclose{\left#1#2\mathclose{}\mathpunct{};#3\right#4}
}
\newcommand{\intff}[2]{\genericinterval{[}{#1}{#2}{]}}
\newcommand{\intoo}[2]{\genericinterval{]}{#1}{#2}{[}}
\newcommand{\intof}[2]{\genericinterval{]}{#1}{#2}{]}}
\newcommand{\intfo}[2]{\genericinterval{[}{#1}{#2}{[}}
\newcommand{\intiff}[2]{\genericinterval{[\![}{#1}{#2}{]\!]}}


%
% Opérateurs mathématiques
%
\DeclareMathOperator{\card}{Card}
\DeclareMathOperator{\tr}{Tr}
\DeclareMathOperator{\Ker}{Ker}
\DeclareMathOperator{\Vect}{Vect}
\DeclareMathOperator{\Span}{Span} % Version anglaise de Vect
\DeclareMathOperator{\indic}{\mathbf{1}}
\DeclareMathOperator{\argth}{argth}
\DeclareMathOperator{\Id}{Id}
\DeclareMathOperator{\Gram}{Gram}
\DeclareMathOperator{\diag}{diag}
\DeclareMathOperator{\Mat}{Mat}
\DeclareMathOperator{\Sp}{Sp}
\DeclareMathOperator{\im}{Im}
\DeclareMathOperator{\Cov}{Cov}
\DeclareMathOperator{\Var}{Var}
\DeclareMathOperator{\Supp}{Supp}
\DeclareMathOperator{\sinC}{sinC}
\DeclareMathOperator{\sgn}{sgn}
\DeclareMathOperator{\argmin}{arg\, min}
\DeclareMathOperator{\cond}{cond} % Fonction conditionnement en optimisation linéaire
\DeclareMathOperator{\proj}{proj}
\DeclareMathOperator{\Bias}{Bias}
\DeclareMathOperator{\MSE}{MSE}
\DeclareMathOperator{\pgcd}{pgcd}
\DeclareMathOperator{\ppcm}{ppcm}
\DeclareMathOperator{\KL}{KL} % Kullback-Leibler divergence


% Pour des symboles « inférieur ou égal », « supérieur ou égal », « ensemble vide » et « parallèles » conformes aux usages français.
\DeclareSymbolFont{AmsA}{U}{msa}{m}{n}
\SetSymbolFont{AmsA}{bold}{U}{msa}{b}{n}
\DeclareMathSymbol\leq\mathrel{AmsA}{"36}
\DeclareMathSymbol\geq\mathrel{AmsA}{"3E}
\DeclareSymbolFont{AmsB}{U}{msb}{m}{n}
\SetSymbolFont{AmsB}{bold}{U}{msb}{b}{n}
\DeclareMathSymbol\emptyset\mathord{AmsB}{"3F}
\def\parallel{\mathrel{/\!/}}


\DeclareMathOperator{\pere}{Père}
\DeclareMathOperator{\DFS}{DFS}
\DeclareMathOperator{\val}{val}

\renewcommand*\rmdefault{ppl}

\begin{document}

\maketitle

\section{Algorithme et complexité}

	\begin{defn}
	\textbf{Algorithme} : suite finie d'instruction non ambiguës pouvant être exécutées de façon automatique.
	\end{defn}

	\begin{defn}
	\textbf{Complexité} de $\mathcal{A}$ : ordre de grandeur du nombre d'opérations élémentaires effectuées pendant le déroulement de l’algorithme.
	\end{defn}

\section{Premières structures de données}

	\begin{defn}
	\textbf{Structure de données} : manière d’organiser et de représenter des données ainsi que méthodes d’accès et de transformation.
	\end{defn}

	\begin{defn}
	\textbf{Liste} : suite ordonnée d’éléments d’un type donné.
	Implémentations possibles : tableau, liste chaînée.
	\end{defn}

	\begin{defn}
	\textbf{Pile} (LIFO) : liste dans laquelle l’insertion ou la suppression d’un élément s’effectue toujours à partir de la même extrémité (début).
	Fonctions standards : empiler, dépiler, tester si vide, renvoyer premier élément.
	\end{defn}

	\begin{defn}
	\textbf{File} (FIFO) : liste dans laquelle les insertions s’effectuent d’un même côté (fin) et les suppressions à partir de l’autre extrémité (début).
	Opérations standards : enfiler, défiler.
	\end{defn}

	\begin{defn}
	Un \textbf{arbre binaire} est soit vide, soit constitué d’un nœud racine et de deux sous-arbres binaires disjoints appelés sous-arbre gauche et sous-arbre droit.
	Il est dit \textbf{localement complet} si tout nœud interne a exactement deux fils.
	Il est dit \textbf{parfait} ou presque-complet si, avec $h$ la hauteur de l’arbre, les niveaux de profondeur $p < h$ sont complètement remplis alors que le niveau de profondeur $h$ est rempli en partant de la gauche.
	Il est dit \textbf{équilibré} si, pour tout nœud, les sous-arbres gauche et droit ont des hauteurs qui diffèrent au plus de 1.
	\end{defn}


\section{Recherche et tri}

	\begin{algorithm}[h]
	\caption{\textcolor{RoyalBlue}{Recherche dichotomique} $O(\ln(n))$}
	\Entree{Tableau $T[1,n]$ trié et élément $x$.}
	$\texttt{gauche} \lar 1$, $\texttt{droite} \lar n$ \;
	$\texttt{found} \lar \texttt{false}$ \;
	\Tq{$\neg \texttt{found}$ et $\texttt{gauche} \leq \texttt{droite}$}
	{
		$\texttt{milieu} \lar (\texttt{gauche} + \texttt{droite})/2$ \;
		\Si{$x < T[\texttt{milieu}]$}
		{
			$\texttt{droite} \lar \texttt{milieu} - 1$ \;
		}
		\SinonSi{$x > T[\texttt{milieu}]$}
		{
			$\texttt{gauche} \lar \texttt{milieu} + 1$ \;
		}
		\Sinon{$\texttt{found} \lar \texttt{true}$}
	}
	\Sortie{$\texttt{found}$}
	\end{algorithm}

	\begin{thm}
	Tout algorithme de tri comparatif (fondé sur des comparaisons entre les éléments pour déterminer la permutation correspondant à l'ordre croissant des données) possède une complexité $\Omega(n \log_2(n))$ (au pire et en moyenne).
	\end{thm}

	\begin{algorithm}
	\caption{\textcolor{RoyalBlue}{Tri sélection} ($C \in O(n^2)$)}
	\Entree{Tableau $T[1,n]$.}
	\SetKwFunction{echanger}{echanger}
	\PourTous{$i \in \intiff{1}{n-1}$}
	{
		$i_{\min} \lar i$, $\texttt{min} \lar T[i]$ \;
		\PourTous{$j \in \intiff{i + 1}{n}$}
		{
			\Si{$T[j] < \texttt{min}$}
			{
				$i_{\min} \lar j$, $\texttt{min} \lar T[j]$ \;
			}
		}
		$\echanger(T,i,i_{\min})$ \;
	}
	\end{algorithm}


	\begin{algorithm}
	\caption{\textcolor{RoyalBlue}{Tri insertion} ($C \in O(n^2)$)}
	\Entree{Tableau $T[1,n]$.}
	\PourTous{$i \in \intiff{2}{n}$}
	{
		$\texttt{j} \lar i$, $\texttt{key} \lar T[j]$ \;
		\Tq{$j \geq 2$ et $T[\texttt{j} - 1] > \texttt{key}$}
		{
			$T[j] \lar T[j - 1]$ \;
			$j \lar j - 1$ \;
		}
		$T[j] \lar \texttt{key}$ \;
	}
	\end{algorithm}


	\begin{algorithm}
	\caption{\textcolor{RoyalBlue}{Tri rapide}}
	\Entree{Tableau $T[1,n]$.}
	\SetKwFunction{echanger}{echanger}
	\SetKwFunction{partition}{partition}
	\SetKwProg{func}{Fonction}{}{}
	\func{\partition{$g$,$d$}}
	{
		$p \lar T[g]$, $i \lar g + 1$, $j \lar d$ \;
		\Tq{$i \leq j$}
		{
			\lTq{$i \leq j$ et $T[i] \leq p$}{$i \lar i + 1$}
			\lTq{$T[j] > p$}{$j \lar j - 1$}
			\Si{$i < j$}
			{
				$\echanger{T,i,j}$\;
				$i \lar i + 1$, $j \lar j - 1$\;
			}
		}
		$\echanger{T,g,j}$\;
		\Sortie{$j$}
	}
	\SetKwFunction{triRapide}{triRapide}
	\SetKwProg{proc}{Procedure}{}{}
	\proc{\triRapide{$g$,$d$}}
	{
		\Si{$g < d$}
		{
			$j \lar \partition{g,d}$ \;
			\triRapide{$g$,$j - 1$}\;
			\triRapide{$j + 1$,$d$}\;
		}
	}
	\triRapide{$1$,$n$}\;
	\end{algorithm}

	\begin{defn}
	Un arbre binaire est un \textbf{arbre binaire de recherche} (ABR) s'il est vide ou égal à $(F_g,c,F_d)$ où : $F_g$ et $F_d$ sont des ABRs, toute clé de $F_g$ est inférieure à $c$ et toute clé de $F_d$ est supérieure à $c$.
	\end{defn}

	Un parcours en ordre infixe d'un ABR donne donc la liste triée de ses clés.

	\begin{algorithm}[h]
	\caption{\textcolor{RoyalBlue}{Insertion dans un ABR}}
	\Entree{Racine $r$ de l'arbre et clé $c$ à insérer.}
	\SetKwFunction{inserer}{inserer}
	\SetKwProg{func}{Fonction}{}{}
	\func{\inserer{$r$,$c$}}
	{
		\Si{$r = \texttt{nil}$}
		{
			Soit $n$ un nouveau nœud \;
			$n.\texttt{data} \lar c$, $n.g \lar \texttt{nil}$, $n.d \lar \texttt{nil}$ \;
			\Sortie{n}
		}
		\Sinon
		{
			\lSi{$c \leq r.\texttt{data}$}{$r.g \lar \inserer(r.g, c)$}
			\lSinon{$r.d \lar \inserer(r.d, c)$}
			\Sortie{r}
		}
	}
	\end{algorithm}

	\begin{pop}
	Le tri rapide et le tri par ABR ont une complexité au pire en $O \left( n^2 \right)$ et en moyene en $O(n\ln(n))$.
	\end{pop}

	\begin{defn}
	\textbf{Tas-max} : arbre binaire parfait sur un ensemble totalement ordonné tel que l'étiquette de chaque nœud autre que la racine est inférieure ou égale à l'étiquette de son père.
	\end{defn}

	\begin{algorithm}[h]
	\caption{\textcolor{RoyalBlue}{Montée, descente et tri dans un tas}}
	\SetKwProg{proc}{Procedure}{}{}
	\Entree{Tableau $T[1,n]$}

	\SetKwFunction{montee}{montee}
	\proc{\montee{p}}
	{
		\tcp{$p$ est le numéro du dernier élément.}
		$i \lar p$ \;
		$\texttt{cle} \lar T[p]$ \;
		\Tq{$i \geq 2$ et $\texttt{cle} > T \left[ \left\lfloor \frac{i}{2} \right\rfloor \right]$}
		{
			$T[i] \lar T \left[ \left\lfloor \frac{i}{2} \right\rfloor \right]$ \;
			$i \lar \left\lfloor \frac{i}{2} \right\rfloor$ \;
		}
		$T[i] \lar \texttt{cle}$ \;
	}

	\SetKwFunction{echanger}{echanger}
	\SetKwFunction{descente}{descente}
	\proc{\descente{q,p}}
	{
		\tcp{$q$ est le numéro de l'élément à descendre.}
		$\texttt{found} \lar \texttt{false}$ \;
		$i \lar q$ \;
		$\texttt{cle} \lar T[q]$ \;
		\Tq{$\neg \texttt{found}$ et $2i \leq p$}
		{
			\lSi{$2i = p$}{$i_{\max} \lar p$}
			\Sinon
			{
				\lSi{$T[2i] \geq T[2i + 1]$}{$i_{\max} \lar 2i$}
				\lSinon{$T[2i] \geq T[2i + 1]$}{$i_{\max} \lar 2i + 1$}
			}
			\Si{$\texttt{cle} < T[i_{\max}]$}
			{
				$T[i] \lar T[i_{\max}]$ \;
				$i \lar i_{\max}$ \;
			}
			\lSinon{$\texttt{found} \lar \texttt{true}$}
		}
		$T[i] \lar \texttt{cle}$ \;
	}

	\SetKwFunction{triTas}{triTas}
	\proc{\triTas{p}}
	{
		\lPour{$p$ qui varie de $2$ à $n$}{\montee{p}}
		\Pour{$p$ qui varie de $n$ à $2$}
		{
			$\echanger{T,1,p}$ \;
			$\descente{1,p - 1}$ \;
		}
	}
	\end{algorithm}

	\begin{pop}
	Le tri tas a une complexité en $O(n\ln(n))$.
	\end{pop}


\section{Le hachage}

	\begin{defn}
	\textbf{Fonction de hachage} : application $h \colon \K \mapsto \intiff{0}{m - 1}$ où $\K$ est l'ensemble des clés possibles et $\intiff{0}{m - 1}$ l'ensemble des indices (adresses de hachage).
	\end{defn}

	\begin{defn}
	\textbf{Collision} primaire : deux clés ont le même indice.
	Collision secondaire : une case d'indice $i$ est déjà occupée par une clé d'indice $j \neq i$.
	\end{defn}

	\begin{defn}
	Hachage linéaire : en cas de collision on se déplace cycliquement vers la droite dans la table.
	Hachage avec chaînage interne : en cas de collision on remplit par la droite une zone de débordement à la fin du tableau, puis les positions libres du tableau, par la droite encore.
	Hachage avec chaînage externe : liste chaînée à partir de chaque case du tableau.
	\end{defn}


\section{L'algorithme de Huffman}

	\begin{defn}
	Codage binaire : application injective $E \colon \Sigma \to \{ 0,1 \}^+$ avec $\Sigma$ un alphabet.
	\end{defn}

	On peut l'étendre à $\Sigma^+$ par concaténation et on dit que $E$ est uniquement décodable si $E$ est injectif sur $\Sigma^+$.

	\begin{defn}
	\textbf{Codage préfixe} : $\forall \sigma, \sigma' \in \Sigma, \sigma \neq \sigma'$, $E(\sigma)$ n'est pas un préfixe de $E(\sigma')$ (suit la règle du préfixe).
	\end{defn}

	\begin{pop}
	Un codage préfixe est uniquement décodable et équivalent à la donnée d'un arbre binaire dont les feuilles sont étiquetées par les lettres de $\Sigma$.
	\end{pop}

	Objectif : associer à la suite des caractères à coder un AB tel que, toute feuille étant associé à un caractère $c_i$ de fréquence $f_i$, à la hauteur $l_i$, $\sum_i l_i f_i$ est minimum.
	L'arbre est alors dit optimum.

	\begin{lem}
	Dans un AB optimum tout nœud interne a deux fils et les deux plus petites occurrences se trouvent à la profondeur maximum de l'arbre.
	\end{lem}

	\begin{lem}
	Il existe un AB optimum dans lequel les deux plus petites occurrences sont dans des feuilles “frères” à la profondeur maximum.
	\end{lem}

	\begin{lem}
	Soit $A$ un AB où les deux plus petites occurrences sont “frères” à la profondeur maximum.
	On note $x$ et $y$ leurs caractères contenus.
	On transforme $A$ en $A'$ en fusionnant ces deux feuilles en $x + y$ avec nombre d'occurrences conservé.
	L'arbre $A$ est optimum si et seulement si $A'$ l'est.
	\end{lem}

	\begin{thm}
	La compression optimale de données sans perte est toujours possible en utilisant un code sans préfixe.
	\end{thm}


	\begin{algorithm}
	\caption{\textcolor{RoyalBlue}{Algorithme de Huffman}}
	\SetKwFunction{push}{push}
	\SetKwFunction{occ}{occ}
	\SetKwFunction{size}{size}
	\SetKwFunction{popBack}{popBack}
	\SetKwFunction{sort}{sort}
	\Entree{alpabet $\Sigma$ et fonction occurrence $\occ$}
	soit $A$ un tableau dynamique d'arbres \;
	\PourTous{$x \in \Sigma$}
	{
		$A.\push{\texttt{nil},(x,\occ(x)),\texttt{nil}}$\;
	}
	$A.\sort{}$ \;
	\tcp{tri par occurrence décroissante des racines}
	$n \lar A.\size{}$ \;
	\Tq{$n > 1$}
	{
		$(\_,(x,\_),\_) \lar A[n]$ \;
		$(\_,(y,\_),\_) \lar A[n - 1]$ \;
		$A[n - 1] \lar (A[n - 1],(x + y, \occ(x) + \occ(y)), A[n])$ \;
		$A.\popBack{}$, $n \lar n - 1$ \;
		$A.\sort{}$ \;
	}
	\Sortie{$A[1]$}
	\end{algorithm}

	\begin{thm}
	L'algorithme de Huffman permet de construire un codage préfixe optimal.
	\end{thm}


\section{Graphes et arbres couvrant}

	On note $X_2 = \{ Y \in \parties(X) \mid \card(Y) = 2 \}$.

	\begin{defn}
	Graphe (simple) non orienté : $G = (V,E)$ avec $E \subset V_2$.
	\end{defn}

	\begin{defn}
	Graphe (simple) orienté : $G = (V,E)$ avec $E \subset V^2$.
	\end{defn}

	\begin{defn}
	L'\textbf{ordre} d'un graphe est $\abs{V}$ et sa \textbf{taille} est $\abs{E}$.
	\end{defn}

	\begin{defn}
	Graphe \textbf{complet} (ou clique) $K_n = (V,V_2)$ avec $\abs{V} = n$.
	\end{defn}

	\begin{defn}
	Graphe \textbf{partiel} de $G = (V,E)$ : $G' = (V,E')$ avec $E' \subset E$.
	Soit $W \subset V$, on appelle $F = (W,E_W)$, avec $E_W = E \cap W_2$ ou $E_W = E \cap W^2$, sous-graphe de $G$ engendré par $W$.
	\end{defn}

	\begin{defn}
	\textbf{Arbre} : graphe acyclique connexe.
	\end{defn}

	\begin{thm}
	Pour un graphe $G$ d'ordre $n$ et de taille $m$, sont équivalents :
		\begin{enumerate}[i)]
		\item $G$ est un arbre,
		\item $G$ est connexe et on a $m = n - 1$,
		\item $G$ est sans cycle et on a $m = n - 1$,
		\item $G$ est connexe et supprimer une arête le déconnecte,
		\item $G$ est acyclique et l'ajout d'une arête crée un cycle,
		\item entre deux sommets quelconques, il existe une chaîne élémentaire unique.
		\end{enumerate}
	\end{thm}

	Problème : étant donné un graphe $G = (V,E)$ pondéré par $w \colon E \to \R$, trouver un arbre couvrant $T$ (graphe partiel connexe) de poids $w(T) = \sum_{e \in T} w(e)$ minimal.

	\begin{algorithm}
	\caption{\textcolor{RoyalBlue}{Algorithme de Kruskal} $O(m \log(m))$}
	\SetKwFunction{triCroissant}{triCroissant}
	\Entree{graphe connexe $G = (V,E)$ et fonction coût $w$}
	$A \lar \emptyset$\;
	$E_t = \triCroissant{E}$ \;
	\PourTous{$(a,b) \in E_t$}
	{
		\Si{$A \cup \{ (a,b) \}$ est acyclique}
		{
			$A \lar A \cup \{ (a,b) \}$ \;
		}
	}
	\Sortie{$(V,A)$}
	\tcc{on peut vérifier l'acyclicité par un tableau donnant le numéro de composante connexe ou une structure union-find}
	\end{algorithm}


	\begin{algorithm}
	\caption{\textcolor{RoyalBlue}{Algorithme de Prim} $O(m \log(n))$}
	\SetKwFunction{voisins}{voisins}
	\Entree{graphe connexe $G = (V,E)$ et fonction coût $w$}
	$S \lar \{ s_o \}$\; \tcp{sommet choisi arbitrairement}
	$A \lar \emptyset$\;
	$p \lar s_o$ \;
	\lPourTous{$s \in V \setminus \{ s_0 \}$}{$d(x) \lar +\infty$}
	\Tq{$S \neq V$}
	{
		\PourTous{$s \in \voisins{p} \cap \compl{S}$}
		{
			\Si{$w(p,s) < d(s)$}
			{
				$\texttt{proche}[s] \lar p$\;
				$d(s) \lar w(p,s)$ \;
			}
		}
		$p \lar \argmin_{s \in \compl{S}}(d(s))$\;
		$S \lar S \cup \{ s \}$ \;
		$A \lar A \cup \{ (\texttt{proche}[p],p) \}$ \;
	}
	\Sortie{$(S,A)$}
	\tcp{on représente $\compl{S}$ avec un tas}
	\end{algorithm}


\section{Problèmes de plus court chemin}

	\begin{note}
	On note $\sim$ la relation binaire telle que $x \sim y$ s'il existe un chemin de $x$ à $y$ dans le graphe.
	\end{note}

	\begin{defn}
	\textbf{Chemin élémentaire} : passant par des sommets distincts.
	\textbf{Circuit} : chemin dont les extrémités coïncident.
	Circuit \textbf{absorbant} $C$ : tel que $\sum_{e \in C} w(e) < 0$.
	\end{defn}

	\begin{defn}
	\textbf{Racine} de $G$ : $r \in V$ tel que $\forall s \in V, r \sim s$.
	\textbf{Arborescence} de racine $r$ : graphe orienté tel que le graphe orienté sous-jacent $G'$ est un arbre et $\forall v \in V$, l'unique chaîne entre $r$ et $v$ dans $G'$ correspond à un chemin de $r$ vers $x$ dans l’arborescence.
	\end{defn}


	\begin{algorithm}
	\caption{\textcolor{RoyalBlue}{Algorithme de Dijkstra}, calcul des plus courts chemins d'un sommet à tous les autres, $O(n^2)$}
	\Entree{graphe $G = (V,E)$ avec coût $w \colon E \to \R_+$ et sommet $s \in V$.}
	$S \lar \{ s \}$\;
	\PourTous{$v \in V$}
	{
		$d[v] \lar w(s,v)$
	}
	\Tq{$\compl{S} \neq \emptyset$}
	{
		$u \lar \argmin_{v \in \compl{S}}(d[v])$\;
		$S \lar S \cup \{ u \}$\;
		\PourTous{$v \in \compl{S}$}
		{
			\Si{$d[v] > d[u] + w(u,v)$}
			{
				$d[v] \lar d[u] + w(u,v)$ \;
				$\texttt{pere}[v] \lar u$ \;
			}
		}
	}
	\Sortie{vecteurs $d$ (distances à $s$) et $\texttt{pere}$}
	\end{algorithm}

	\begin{defn}
	On appelle \textbf{tri topologique} des sommets d'un graphe orienté acyclique une numérotation des sommets telle que, pour $(u,v) \in E$ et $n$ la fonction de numérotation, $n(u) \leq n(v)$.
	\end{defn}

	\begin{pop}
	Un graphe est sans circuit si et seulement s'il admet une numérotation topologique.
	\end{pop}

	\begin{algorithm}
	\caption{\textcolor{RoyalBlue}{Algorithme de Bellman} $O(n^2)$}
	\Entree{graphe acyclique orienté avec coût w $G = (V,E)$}
	\SetKwFunction{numT}{numTopo}
	\SetKwProg{func}{Fonction}{}{}
	\func{\numT{$V$,$E$}}
	{
		trouver $x$ qui n'a pas de prédecesseur \;
		$n_{-1} = \numT(V \setminus \{ x \}, E \cap (V \setminus \{ x \})^2)$ \;
		\Sortie{$n \colon y \mapsto \left\{ \begin{array}{cl}	1 & \text{si } y = x \\ n_{-1}(y) + 1 & \text{sinon} \end{array}\right.$}
	}
	$n = \numT{V,E}$ \;
	$r \lar n^{-1}(1)$ \;
	$d[r] \lar 0$ \;
	\lPourTous{$v \in V \setminus \{ r \}$}{$d[v] \lar \infty$}
	\Pour{$i$ variant de $2$ à $n$}
	{
		$x \lar n^{-1}(i)$ \;
		$\texttt{pere}[x] \lar \argmin_{y \in V \mid n(y) < y, (y,x) \in E}(d[y] + w(y,x))$ \;
		$d[x] \lar d[\texttt{pere}[x]] + w(\texttt{pere}[x],x)$ \;
	}
	\Sortie{vecteurs $d$ et $\texttt{pere}$}
	\end{algorithm}


	\begin{algorithm}
	\caption{\textcolor{RoyalBlue}{Algorithme de Ford}}
	\Entree{Graphe $G = (V,E)$ avec coût $w \colon E \to \R$ et sommet $s \in V$.}
	$k \lar 0$\;
	$d^{(0)} (s) \lar 0$\;
	\lPourTous{$v \in V \setminus \{ s \}$}{$d^{(0)}(v) \lar \infty$}
	\Repeter{$k = n$ ou $\neg \texttt{changement}$}
	{
		$\texttt{changement} \lar \texttt{faux}$ \;
		$k \lar k+1$ \;
		\PourTous{$v \in V$}
		{
			$d^{(k)}(v) \lar \min_{u \mid (u,v) \in E} (d^{(k-1)} (u) + w(u,v))$ \;
			\Si{$d^{(k)} (v) \neq d^{(k-1)} (v)$}
			{
				$\texttt{changement} \lar \texttt{vrai}$ \;
				$\texttt{pere}(v) \lar \underset{u \mid (u,v) \in E}{\argmin} (d^{(k-1)} (u) + w(u,v))$ \;
			}
		}
	}
	\Si{$\texttt{changement}$}
	{
		\Sortie{"cycle absorbant"}
	}
	\Sinon
	{
		\Sortie{vecteurs $d$ et $\texttt{pere}$}
	}
	\end{algorithm}


	\begin{algorithm}
	\caption{\textcolor{RoyalBlue}{Algorithme de Dantzig}, plus courts chemins de tout sommet à tout sommet, cas général $O(n^3)$}
	\Entree{graphe $G = (V,E)$ avec coût $w \colon E \to \R$}
	$D^{(1)}(1,1) \lar 0$\;
	\Pour{$k$ allant de $1$ à $n - 1$}
	{
		\Pour{$i \in \intiff{1}{k}$}
		{
			$D^{(k + 1)}(i,k + 1) \lar \underset{j \in \intiff{1}{k}}{\min}(D^{(k)}(i,j) + w(j,k + 1))$\;
			$D^{(k + 1)}(k + 1,i) \lar \underset{j \in \intiff{1}{k}}{\min}(D^{(k)}(j,i) + w(k + 1,j))$\;
		}
		$D^{(k + 1)}(k + 1,k + 1) \lar 0$ \;
		\Pour{$(i,j) \in \intiff{1}{k}^2$}
		{
			$D^{(k + 1)}(i,j) \lar \min(D^{(k)}(i,j), D^{(k + 1)}(i,k + 1) + D^{(k + 1)}(k + 1,j))$\;
		}
	}
	\Sortie{matrice $D^{(n)}$ des plus courtes distances}
	\end{algorithm}



\section{Parcours de graphes}

	\begin{algorithm}
	\caption{\textcolor{RoyalBlue}{Parcours abstrait d'un graphe}}
	\Entree{$G = (V,E)$, sommet $r \in V$.}
	$T \leftarrow \{ r \}$\;
	marquer $r$ \;
	\Tq{$T \neq \emptyset$}
	{
		choisir $v \in T$\;
		visiter $v$ \;
		$T \leftarrow T \setminus \{ v \}$ \;
		\PourTous{$v'$ voisin de $v$}
		{
			\Si{$v'$ n'est pas marqué}
			{
				$T \lar T \cup \{ v' \}$ \;
				marquer $v'$ \;
			}
		}
	}
	\end{algorithm}

	C'est sur la façon de choisir $v$ que diffèrent les algorithmes : parcours en largeur en utilisant une file et parcours en profondeur en utilisant une pile.
	
	\begin{note}
	$M$ : ensemble des sommets marqués par un algorithme de parcours à partir de $r$.
	\end{note}

	\begin{defn}
	\textbf{Arborescence du parcours} : $A = (M, \{ (\texttt{pere}(x),x) \mid x \in M \setminus \{ r \}\})$ de racine $r$.
	Les arcs $(\texttt{pere}(x),x)$ s'appellent \textbf{arcs arborescents}.
	\end{defn}

	\begin{pop}
	$M$ est l'ensemble des sommets $x$ pour lesquels il existe dans $G$ un chemin de $r$ à $x$.
	\end{pop}

	\begin{algorithm}
	\caption{\textcolor{RoyalBlue}{Parcours en profondeur} $O(n + m)$}
	\SetKwProg{func}{Fonction}{}{}
	\Entree{$G = \{ V,E \}$, $r$ un sommet particulier.}
	($n_{pre} \lar 1$, $n_{post} \lar 1$) \;
	\SetKwFunction{DFS}{DFS}
	\func{\DFS{$s \in V$}}
	{
		(Numéroter $s$ en préfixe, $n_{pre} \lar n_{pre} + 1$) \;
		Marquer $s$ \;
		\PourTous{$v$ voisin de $s$}
		{
			\Si{$v$ n'est pas marqué}
			{
				$\texttt{pere}(v) \lar s$ \;
				$\DFS(v)$ \;
			}
		}
		(Numéroter $s$ en suffixe, $n_{post} \lar n_{post} + 1$) \;
	}
	\DFS{$r$}\;
	\end{algorithm}

	Dans un DFS en non orienté il peut être utile d'orienter les arêtes traversées (lors de la découverte des voisins) dans le sens de cette traversée.

	\begin{defn}
	\textbf{Arête arborescente} : arête qui a été orientée en un arc arborescent.
	\textbf{Arc arrière} : arc obtenu en orientant une arête non arborescente.
	Toute arête conduit à un arc arborescent ou à un arc arrière.
	\end{defn}

	\begin{lem}
	Les arcs arrières vont d'un sommet $x$ à un des ses ancêtres autre que son père dans l'arborescence.
	\end{lem}

	\begin{defn}
	Soit $\Join$ la relation d'équivalence sur $V$ définie par $x \Join y \iff (x \sim y) \wedge (y \sim x)$.
	Les classes d'équivalences de $x$ pour cette relation sont les composantes fortement connexes de $G$.
	$G$ est dit fortement connexe s'il n'admet qu'une seule telle composante.
	\end{defn}

	\begin{algorithm}
	\caption{\textcolor{RoyalBlue}{Algorithme de calcul des composantes fortement connexes} $O(n + m)$}
	\Entree{$G = \{ V,E \}$ un graphe orienté.}
	\SetKwFunction{DFS}{DFS}
	%$n_{post} \colon \intiff{1}{n} \to \intiff{1}{n}$ \;
	\Tq{tous les sommet ne sont pas numérotés en post-fixe}
	{
		soit $x$ un sommet non numéroté \;
		$\DFS(x)$ avec numérotation post-fixe \;
	}
	$E^- \lar \{ (v,u) \mid (u,v) \in E \}$ \;
	$G^- \lar \{ V,E^- \}$ \;
	effacer les marquages \;
	effectuer des DFS dans $G^-$ successivement au départ du sommet de plus grand numéro post-fixe restant
	\tcp{chaque DFS donne une composante fortement connexe}
	\Sortie{composantes fortement connexes}
	\end{algorithm}

	\begin{defn}
	\textbf{Sommet d'articulation} d'un graphe non orienté : sommet dont la suppression augmente le nombre de composantes connexes.
	C'est équivalent à dire qu'il existe $(x,y) \in V^2$ tel que toute chaîne entre $x$ et $y$ passe par ce point.
	Un graphe est dit \textbf{$2$-connexe} si et seulement si il est réduit à une arête ou n'a pas de sommet d'articulation.
	\end{defn}

	\begin{thm}
	Soit $G = (V,E)$ un graphe connexe, non orienté, et $A$ une arborescence DFS de $G$.
	Alors $s \in V$ est un point d'articulation de $G$ si et seulement si
		\begin{itemize}
		\item $s$ est la racine de $A$ et possède au moins deux fils dans $A$,
		\item ou $s$ n'est pas la racine de $A$ et, pour un fils $f$ de $s$ dans $A$, il n'y a pas d'arc arrière entre les descendants de $f$ ($f$ compris) et les ancêtres de $s$.
		\end{itemize}
	\end{thm}

	\begin{defn}
	\textbf{Numérotation basse} : $\texttt{basse}(v)$ est le minimum de : $\texttt{prefixe}(v)$, $\texttt{prefixe}(z)$ où $z \in V$ et $(v,z)$, et $\texttt{basse}(u)$ où $u$ est fils de $v$ dans l'arborescence DFS.
	\end{defn}

	\begin{thm}
	Le point de départ de DFS est sommet d'articulation si et seulement si il a au moins deux fils dans l'arborescence.
	Un sommet $s$ autre que la racine est un sommet d'articulation si et seulement s'il a au moins un fils $f$ dans l'arborescence tel que $\texttt{basse}(f) \geq \texttt{prefixe}(s)$.
	\end{thm}

	\begin{defn}
	Soit $G$ un graphe non orienté.
	On dit que $G$ est $2$-connexe s'il est connexe et n'admet aucun sommet d'articulation.
	On appelle composante $2$-connexe tout sous-graphe $2$-connexe maximal pour l'inclusion.
	\end{defn}

	\begin{algorithm}
	\caption{\textcolor{RoyalBlue}{Calcul des composantes $2$-connexes} $O(m)$}
	\SetKwProg{func}{Fonction}{}{}
	\Entree{$G = \{ V,E \}$, $r$ un sommet particulier.}
	$\texttt{pre} \lar 1$ \;
	$\texttt{pileAretes} \lar \texttt{nil}$ \;
	\SetKwFunction{empiler}{empiler}
	\SetKwFunction{parcours}{parcours}
	\func{\parcours{$x \in V$}}
	{
		Marquer $x$ \;
		$\texttt{prefixe}[x] \lar \texttt{pre}$ \;
		$\texttt{pre} \lar \texttt{pre} + 1$ \;
		$\texttt{basse}[x] \lar \texttt{prefixe}[x]$ \;
		\PourTous{$y$ voisin de $x$}
		{
			\Si{$\{ x,y \}$ n'a pas été traversée}
			{
				\tcc{i.e. si $y$ n'est pas marqué ou ($y$ n'est pas le père de $x$ et $\texttt{prefixe}[y] \leq \texttt{prefixe}[x]$)}
				$\texttt{pileAretes}.\empiler{ \{ x,y \} }$
				\Si{$y$ est déjà marqué}
				{
					$\texttt{basse}[x] \lar \min(\texttt{basse}[x], \texttt{prefixe}[y])$ \;
				}
				\Sinon
				{
					$\texttt{pere}[y] \lar x$ \;
					\parcours{$y$} \;
					\Si{$\texttt{basse}[y] \geq \texttt{prefixe}[x]$}
					{
						dépiler $\texttt{pileAretes}$ jusqu'à $\{ x,y \}$ (compris) \;
						les arêtes dépilées sont celles d'une composante $2$-connexe \;
					}
					\Sinon
					{
						$\texttt{basse}[x] \lar \min(\texttt{basse}[x], \texttt{basse}[y])$ \;
					}
				}
			}
		}
	}
	\parcours{$r$}\;
	\end{algorithm}


\section{Flot max et coupe min}

	\begin{defn}
	\textbf{Réseau} : quadruplet $(G,c,s,t)$ où $G = (V,E)$ est un graphe orienté, $c \colon E \to \R_+^*$ une fonction de capacité et $s$ et $v$ des sommets distincts de $V$ appelés source et puits.
	\end{defn}

	\begin{note}
	Soit $S \subset V$ contenant $s$ et non $p$.
	On dit que $S$ sépare $s$ de $p$.
	\textbf{Coupe} : ensemble noté $(S,\compl{S})$ des arcs d'origine dans $S$ et d'extrémité dans $\compl{S}$.
	On note $E(S,T) = \{ (u,v) \in E \mid u \in S, v \in T \}$ et $\delta_f(v) = \sum_{e \in E(v,V \setminus \{ v \})} f(e) - \sum_{e \in E(V \setminus \{ v \},v)} f(e)$.
	\end{note}

	\begin{defn}
	\textbf{Flot}: application $f \colon E \to \R_+$ telle que
		\begin{enumerate}[i)]
		\item $\forall u \in E, f(e) \leq c(e)$
		\item $\forall v \in V \setminus \{ s,t \}, \delta_f(v) = 0$.
		\end{enumerate}
	Flux d'un arc $u$ : $f(u)$.
	\end{defn}

	\begin{defn}
	\textbf{Valeur du flot} $f$ : $\val(f) = \delta_f(s) = - \delta_f(p) = f(S,\compl{S}) - f(\compl{S},S)$ avec $(S,\compl{S})$ une coupe.
	En particulier cette dernière valeur est indépendante de la coupe choisie.
	Capacité d'une coupe : $c(S,\compl{S}) = \sum_{u \in (S,\compl{S})} c(u)$.
	\end{defn}

	\begin{defn}
	Étant donné un réseau, un flot maximal est un flot qui maximise $\val{f}$.
	Problème :  trouver un tel flot.
	\end{defn}

	\begin{thm}
	Soit $f$ un flot et $(S,\bar{S})$ une coupe, alors
		\begin{enumerate}
		\item $\val(f) \leq c(S,\compl{S})$,
		\item si $\val(f) = c(S,\compl{S})$, $\val(f)$ est maximum et $c(S,\compl{S})$ est minimum,
		\item $\val(f) = c(S,\compl{S})$ si et seulement si $\forall u \in (S,\compl{S}), f(u) = c(u)$ et $\forall u \in (\compl{S},S), f(u) = 0$.
		\end{enumerate}
	\end{thm}

	\begin{note}
	Soit $e = (u,v)$ un arc, on note $\overleftarrow{e} = (v,u)$ (arc inverse).
	Pour $f$ donné, on a le graphe résiduel $(G_f,c',s,t)$ où $G_f = (V,E')$, $E' = \{ e \in E \mid f(e) < c(e) \} \cup \{ \overleftarrow{e} \mid e \in E, f(e) > 0 \}$ et $\forall e \in E', c'(e) = \left\{ \begin{array}{cl}
		c(e) - f(e) & \text{ si } e \in E \\
		f(\overleftarrow{e}) & \text{ si } \overleftarrow{e} \in E
	\end{array}\right.$.
	\end{note}

	\begin{defn}
	\textbf{Chemin $f$-augmentant} $P$ : chemin de $s$ à $t$ dans $G_f$.
	\end{defn}

	Soit $\gamma = \min_{e \in P} c'(e)$.
	On peut augmenter $f$ le long de $P$ par les formules $f(e) \lar f(e) + \gamma$ si $e \in E$ et $f(e') \lar f(e') - \gamma$ si $\overleftarrow{e} \in E$.

	\begin{thm}[de Ford-Fulkerson]
	Sont équivalents
		\begin{enumerate}[i.]
		\item $f$ est un flot maximal,
		\item $G_f$ ne contient pas de chemin $f$-augmentant,
		\item $\exists (S,\bar{S})$ une coupe, telle que $\abs{f} = c(S,\bar{S})$.
		\end{enumerate}
	\end{thm}


	\begin{algorithm}
	\caption{\textcolor{RoyalBlue}{Algorithme de Ford-Fulkerson} $O(m \cdot \max_f(\val(f)))$}
	\SetKwProg{func}{Fonction}{}{}
	\Entree{$(G,c,s,t)$}
	\SetKwFunction{rca}{rechercheCheminAugmentant}
	\func{\rca{}}
	{
		$\texttt{marque}[s] \lar (\Delta,+\infty)$ \;
		considérer les sommets de $V \setminus \{ s \}$ non marqués \;
		considérer tous les sommets non examinés \;
		\Tq{$p$ est non marqué et il y a un sommet marqué mais non examiné}
		{
			soit $x$ un tel sommet \;
			\tcp{on prolonge la chaîne augmentant après $x$}
			soit $\alpha$ la valeur absolue de la deuxième marque de $x$\;
			\PourTous{$y$ successeur de $x$ non marqué}
			{
				\Si{$c(x,y) > f(x,y)$}
					{$\texttt{marque}[y] \lar (x,\min(\alpha,c(x,y) - f(x,y)))$ \;}
			}
			\PourTous{$z$ prédécesseur de $x$ non marqué}
			{
				\Si{$f(z,x) > 0$}
					{$\texttt{marque}[z] \lar (x,-\min(\alpha,f(z,x)))$\;}
			}
			considérer $x$ comme examiné \;
		}
	}
	on initialise $f$ à $0$ (flot nul) \;
	\rca{} \;
	\Tq{$p$ est marqué}
	{
		reconstituer la chaîne augmentante $C$ \;
		soit $\alpha$ la valeur absolue de la deuxième marque de $p$\;
		\PourTous{arc $(x,y) \in C$ parcouru à l'endroit}
		{
			$f(x,y) \lar f(x,y) + \alpha$ \;
		}
		\PourTous{arc $(x,y) \in C$ parcouru à l'envers}
		{
			$f(x,y) \lar f(x,y) - \alpha$ \;
		}
		\rca{} \;
	}
	\Sortie{$f$}
	\end{algorithm}

	\begin{thm}
	Dans un réseau à capacités entières, il existe un flot maximum tel que tous les flux soient entiers.
	\end{thm}


\section{Applications de la théorie des flots}

	\begin{defn}
	$G$ est dit \textbf{$k$-connexe} s’il a au moins $k + 1$ sommets et si la suppression d’au plus $k - 1$ sommets quelconques résulte en un graphe connexe.
	Sommet-connectivité de $G$ : plus grand entier $k$ tel que $G$ soit $k$-connexe.
	\end{defn}

	\begin{defn}
	$G$ est dit \textbf{$k$-arête-connexe} si la suppression d’au plus $k - 1$ arêtes quelconques résulte en un graphe connexe.
	Arête-connectivité de $G$ : plus grand entier $k$ tel que $G$ soit $k$-arête-connexe.
	La forte arc-connectivité est analogue dans un graphe orienté.
	\end{defn}

	\subsection{Forte arc-connectivité}

		\begin{note}
		Soit $G = (V,e)$ orienté et $(a,b) \in V^2$ distincts.
		On considère le réseau $R_{ab}$ déterminé par : ce graphe $G$, le sommet source $a$ et le sommet puits $b$, une capacité de $1$ sur tous les arcs de $G$.
		\end{note}

		\begin{note}
		Soit $P$ et $N$ de à valeurs de $E$ dans $\N$ tels que $P(a,b)$ est le nombre maximum de chemins de $a$ vers $b$, deux à deux arc-disjoints, et $N(a,b)$ le nombre minimum d’arcs à supprimer pour qu’il n’existe plus de chemin de $a$ vers $b$.
		\end{note}

		\begin{thm}[Menger, “Max Flow - Min Cut”]
		Si $f_{ab}$ est maximal dans le réseau $R_{ab}$, alors $\val(f_{ab}) = P(a,b) = N(a,b)$.
		\end{thm}

		\begin{lem}[Zorn]
		Supposons définie une numérotation des sommets de $G$ de $0$ à $n - 1$.
		La forte arc-connectivité de $G$ est $\min_{0 \leq i \leq n-1} N(x_i,x_{i+1})$, en posant $x_n = x_0$.
		\end{lem}

		Puisque l'algorithme de Ford-Fulkerson est en $O(nm)$, on peut déterminer une forte arc-connectivité en $O(mn^2)$.

	\subsection{Détermination de l'arête-connectivité}

		\begin{note}
		Symétrisé $G^* = (V,E^*)$ de $G = (V,E)$ : $\forall \{ u,v \} \in E$, $(u,v) \in E^*$ et $(v,u) \in E^*$.
		Par ailleurs, on attribue à chaque arc de $G^*$ une capacité de 1.
		\end{note}

		\begin{note}
		Soit $a$ et $b$ dans $V$, on note :
			\begin{itemize}
			\item $N(a,b)$ le nombre min d’arêtes à supprimer de $G$ pour qu’il n’existe plus de chaîne entre $a$ et $b$,
			\item $N^*(a,b)$ le nombre min d’arcs à supprimer de $G^*$ pour qu’il n’existe plus de chemin de $a$ vers $b$,
			\item $N(a,b)$ le nombre max de chaînes deux à deux arête-disjointes de $G$ entre $a$ et $b$,
			\item $N^*(a,b)$ le nombre max de chemins deux à deux arc-disjoints de $G^*$ de $a$ vers $b$.
			\end{itemize}
		\end{note}

		\begin{thm}
		Si $f_{ab}$ de $a$ à $b$ est maximal dans $G^*$, alors $N(a,b) = N^*(a,b) = P(a,b) = P^*(a,b) = \val(f_{ab})$.
		\end{thm}

	\subsection{Forte connectivité et connectivité}

		\begin{thm}[Menger]
		Un graphe non orienté est $k$-connexe si et seulement si, entre deux sommets quelconques, il existe $k$ chaînes n’ayant en commun que leurs extrémités.
		\end{thm}

		\begin{thm}[Menger]
		Un graphe orienté est $k$-fortement connexe si et seulement si, entre deux sommets quelconques, il existe $k$ chemins de l'un vers l'autre n’ayant en commun que leurs extrémités.
		\end{thm}

	\subsection{Couplage maximum}

		\begin{defn}
		\textbf{Graphe biparti} : il existe une bipartition des sommets $V = V_1 \cup V_2$ telle que $E \subset V_1 \times V_2$.
		\textbf{Couplage} : ensemble d’arêtes non incidentes entre elles.
		\end{defn}

		Problème : étant donné un graphe $G = (V,E)$ biparti, trouver un couplage de cardinal maximal.

		Modélisation : on construit un réseau à partir de $G$ en adjoignant deux sommets $s$ et $p$ et en plaçant des arcs de capacité $1$ :
		\begin{itemize}
			\item entre $s$ et $v$ pour tout $v \in V_1$
			\item entre $v$ et $p$ pour tout $v \in V_2$
			\item entre $u$ et $v$ si $(u,v) \in E$
		\end{itemize}

		\begin{pop}
		Les flots à valeur entière sont en bijection avec les couplages de $G$.
		\end{pop}

		L'algorithme de Ford-Fulkerson permet alors de résoudre le problème d'affectation posé.


\section{Complexité}

	\begin{defn}
	Une machine de Turing déterministe à une bande est formée de :
	\begin{itemize}
	\item un ruban de mémoire infini,
	\item une tête de lecture et d'écriture qui se déplace sur le ruban,
	\item un ensemble d'états avec une table de transition (comme une unité centrale munie d'un programme),
	\item un état particulier qui est l'état initial,
	\item un état final
	\end{itemize}
	\end{defn}

	Un calcul sur une machine de Turing se fait :
	\begin{itemize}
	\item en inscrivant l'entrée sur le ruban,
	\item en appliquant la table de transition,
	\item si on arrive dans l'état final le calcul s'arrête et est accepté.
	\end{itemize}

	Modèle : $(Q,\Gamma,b,\Sigma,\delta,q_0,q_s)$ avec $Q$ l'ensemble des états, $\Gamma$ l'alphabet de travail, $b \in \Sigma$ le symbole blanc, $\Sigma \subset \Gamma$ l'alphabet d'entrée, $\delta \colon Q \times \Gamma \to Q \times \Gamma \times \{ \lar, \rar \}$, $q_0$ l'état initial et $q_s$ l'état final.

	\subsection{Classe de complexité $P$}

		\begin{defn}
		Une fonction $f \colon \N \to \R$ est dite polynomiale en $g \colon \N \to \R$ si $\exists p \in \R[g], \forall n \in \N, f(n) \leq p(g(n))$.
		\end{defn}

		\begin{defn}
		\textbf{Problème de décision} (ou problème de reconnaissance) : admet oui ou non pour réponse.
		Il est \textbf{en temps polynomial} (ou appartient à la classe $P$) s'il existe une machine de Turing permettant de décider ce problème et qui fonctionne en un nombre d'étapes polynomial en la taille de l'instance écrite sur le ruban d'entrée.
		\end{defn}

		\begin{note}
		$A$ est polynomialement réductible à $B$ si toute instance $a$ de $A$ peut être transformée en temps polynomial en une instance $b$ de $B$, de longueur polynomiale en la longueur de $a$, telle que $A$ est solvable sur $a$ si et seulement si $B$ est solvable sur $b$.
		On note $A \leq_P B$.
		\end{note}

	\subsection{Classe de complexité $NP$}

		\begin{defn}
		Un problème est dans la classe $NP$ s'il est possible de vérifier une solution en temps polynomial.
		\end{defn}
		
		\begin{pop}
		Le problème du voyageur de commerce est dans la classe $NP$.
		\end{pop}

		\begin{rem}
		$P \subset NP$ mais on ne sait pas dire si il y a égalité ou non.
		\end{rem}

		\begin{defn}
		On dit que $A$ est $NP$-complet si $A \in NP$ et, pour tout $B \in NP$, $B \leq_P A$.
		\end{defn}

		\begin{pop}
		Si $A$ est $NP$-complet et $C \in NP$ tel que $A \leq_P C$, alors $C$ est aussi $NP$-complet.
		\end{pop}

	\subsection{Problème de satisfiabilité}

		\begin{voc}
			\begin{itemize}
			\item Variables booléennes : $x_1,\ldots,x_n$.
			\item Littéraux : $\lambda_i = x_i$ ou $\lambda_i = \neg x_i$.
			\item Clauses : $C_j = \lambda_i \vee \lambda_2 \vee \ldots \vee \lambda_k$.
			\item Formule : $C_1 \wedge C_2 \wedge \ldots \wedge C_z$.
			\end{itemize}
		\end{voc}

		\begin{defn}
		Problème SAT : étant donné une formule $F$, peut-on trouver des valeurs de $x_1,\ldots,x_n$ qui rendent $F$ vraie ?
		\end{defn}

		\begin{thm}[Cook-Levin]
		Le problème SAT est $NP$-complet.
		\end{thm}

		\begin{defn}
		Problème 3-SAT : étant donné une formule $F$ dont les clauses sont toutes formées de 3 littéraux, peut-on satisfaire $F$ ?
		\end{defn}

		\begin{thm}
		Le problème 3-SAT est $NP$-complet.
		\end{thm}
		
		\begin{defn}
		Si le problème de décision associé à une problème d'optimisation $O$ est $NP$-complet, alors $O$ est lui-même $NP$-difficile.
		\end{defn}
		
		\vspace{3cm}
		\begin{figure}[htp]
		\centering
		\includegraphics[scale=0.90]{Ada-Lovelace.png}
		\caption{Ada Lovelace}
		\label{}
		\end{figure}

\end{document}